\documentclass[10pt,a4paper]{amsart}
\usepackage{amsmath,amssymb,amsthm,graphicx}
% Retained minimal preamble for the Stacks Project chapter sites-cohomology.tex (Cohomology on Sites).
% Verbatim excerpt of the chapter preamble: Xy-pic setup (native preamble lines 14-18) and the native
% theorem declarations with the mathematical macros (native preamble lines 154-217). Copied mechanically;
% nothing added, removed or redefined. The wrapper class and the amsmath/amssymb/amsthm/graphicx base
% packages are supplied by the assembler, so the class/externaldocument part of the native preamble is omitted.
% For commutative diagrams we use Xy-pic
\usepackage[all]{xy}

% We use 2cell for 2-commutative diagrams.
\xyoption{2cell}
% Theorem environments.
%
\theoremstyle{plain}
\newtheorem{theorem}[subsection]{Theorem}
\newtheorem{proposition}[subsection]{Proposition}
\newtheorem{lemma}[subsection]{Lemma}

\theoremstyle{definition}
\newtheorem{definition}[subsection]{Definition}
\newtheorem{example}[subsection]{Example}
\newtheorem{exercise}[subsection]{Exercise}
\newtheorem{situation}[subsection]{Situation}

\theoremstyle{remark}
\newtheorem{remark}[subsection]{Remark}
\newtheorem{remarks}[subsection]{Remarks}

\numberwithin{equation}{subsection}

% Macros
%
\def\lim{\mathop{\mathrm{lim}}\nolimits}
\def\colim{\mathop{\mathrm{colim}}\nolimits}
\def\Spec{\mathop{\mathrm{Spec}}}
\def\Hom{\mathop{\mathrm{Hom}}\nolimits}
\def\Ext{\mathop{\mathrm{Ext}}\nolimits}
\def\SheafHom{\mathop{\mathcal{H}\!\mathit{om}}\nolimits}
\def\SheafExt{\mathop{\mathcal{E}\!\mathit{xt}}\nolimits}
\def\Sch{\mathit{Sch}}
\def\Mor{\mathop{\mathrm{Mor}}\nolimits}
\def\Ob{\mathop{\mathrm{Ob}}\nolimits}
\def\Sh{\mathop{\mathit{Sh}}\nolimits}
\def\NL{\mathop{N\!L}\nolimits}
\def\CH{\mathop{\mathrm{CH}}\nolimits}
\def\proetale{{pro\text{-}\acute{e}tale}}
\def\etale{{\acute{e}tale}}
\def\QCoh{\mathit{QCoh}}
\def\Ker{\mathop{\mathrm{Ker}}}
\def\Im{\mathop{\mathrm{Im}}}
\def\Coker{\mathop{\mathrm{Coker}}}
\def\Coim{\mathop{\mathrm{Coim}}}

% Boxtimes
%
\DeclareMathSymbol{\boxtimes}{\mathbin}{AMSa}{"02}

%
% Macros for moduli stacks/spaces
%
\def\QCohstack{\mathcal{QC}\!\mathit{oh}}
\def\Cohstack{\mathcal{C}\!\mathit{oh}}
\def\Spacesstack{\mathcal{S}\!\mathit{paces}}
\def\Quotfunctor{\mathrm{Quot}}
\def\Hilbfunctor{\mathrm{Hilb}}
\def\Curvesstack{\mathcal{C}\!\mathit{urves}}
\def\Polarizedstack{\mathcal{P}\!\mathit{olarized}}
\def\Complexesstack{\mathcal{C}\!\mathit{omplexes}}
% \Pic is the operator that assigns to X its picard group, usage \Pic(X)
% \Picardstack_{X/B} denotes the Picard stack of X over B
% \Picardfunctor_{X/B} denotes the Picard functor of X over B
\def\Pic{\mathop{\mathrm{Pic}}\nolimits}
\def\Picardstack{\mathcal{P}\!\mathit{ic}}
\def\Picardfunctor{\mathrm{Pic}}
\def\Deformationcategory{\mathcal{D}\!\mathit{ef}}
\usepackage{cleveref}
\crefname{originalreference}{original reference}{original references}
\Crefname{originalreference}{Original reference}{Original references}
\makeatletter
\newlabel{lemma-fibred-category-with-object}{{08P9}{source}{}{sourcepointer-1}{}}
\newlabel{lemma-fibred-category-with-object@cref}{{[originalreference][0][]08P9}{[1][0][]source}{}{}{}}
\makeatother
\crefname{originalreference}{original reference}{original references}
\Crefname{originalreference}{Original reference}{Original references}
\makeatletter
\newlabel{situation-morphism-fibred-categories}{{08PA}{source}{}{sourcepointer-2}{}}
\newlabel{situation-morphism-fibred-categories@cref}{{[originalreference][0][]08PA}{[1][0][]source}{}{}{}}
\makeatother
\title{For a morphism of fibred categories the sections of the localised comparison morphisms are compatible}
\author{The Stacks Project Authors}
\date{stacks-project:sites-cohomology.tex:local-snapshot-20261011; GNU FDL 1.2 or later}
\begin{document}
\maketitle
\noindent\textit{Editorial note.} All delivered mathematical text is the author's own, copied from the pinned Stacks Project chapter file sites-cohomology.tex; no statement or proof step was rewritten, reordered or removed.\par\medskip
\noindent\textit{Editorial note.} The counted proof is the author's complete argument: the proof is written out in full: the construction of the functor v' and its identification with the corresponding functor for p, and the transfer of the identities from the single-category lemma.\par\medskip
\noindent\textit{Editorial note.} The label lemma-morphism-fibred-categories-with-object names the counted statement, and the proof environment immediately following it in the source is delivered verbatim as its proof.\par\medskip
\noindent\textit{Editorial note.} The counted statement is carried by the author's own lemma environment in the chapter source.\par\medskip
\noindent\textit{Editorial note.} The statement's displayed diagram is the author's and is delivered verbatim; the two cited inputs are kept as external pointers with their canonical Stacks tags.\par\medskip
\noindent\textit{Allowed standard prerequisites (editorial).}\begin{itemize}
\item The section of the localised morphism for a single fibred category. Same chapter, the lemma constructing the functor v : D/p(U) -> C/U and the section sigma with pi\_U o sigma = id and sigma\textasciicircum{}\{-1\} = pi\_\{U,*\}. Supplies the construction applied twice, to p and to p'.
\item Morphisms of fibred categories preserve strongly cartesian morphisms. Same chapter, the standing situation of a morphism of fibred categories over the base, whose defining property is the preservation of strongly cartesian morphisms. Supplies that u o v' is the functor attached to p and U.
\end{itemize}\medskip
\ifdefined\hypertarget\hypertarget{sourcepointer-1}{}\fi
\noindent\textit{Source reference.} Original reference 08P9: sources/stacks-sites-cohomology/pinned-original-source.tex, source lines 9720--9720. Prerequisite pointer invoked at proof line(s) 9818; it names a standard fact of the same chapter that the counted argument uses, and every step of the counted proof built on it is retained verbatim in the delivered body.\par\medskip
\ifdefined\hypertarget\hypertarget{sourcepointer-2}{}\fi
\noindent\textit{Source reference.} Original reference 08PA: sources/stacks-sites-cohomology/pinned-original-source.tex, source lines 9767--9767. Prerequisite pointer invoked at proof line(s) 9797; it names a standard fact of the same chapter that the counted argument uses, and every step of the counted proof built on it is retained verbatim in the delivered body.\par\medskip
\begin{lemma}
\label{lemma-morphism-fibred-categories-with-object}
Assumptions and notation as in
Situation \ref{situation-morphism-fibred-categories}.
For $U' \in \Ob(\mathcal{C}')$ set $U = u(U')$ and $V = p'(U')$ and
consider the induced morphisms of ringed topoi
$$
\xymatrix{
(\Sh(\mathcal{C}'/U'), \mathcal{O}_{U'}) \ar[rd]_{\pi'_{U'}} \ar[rr]_{g'} & &
(\Sh(\mathcal{C}), \mathcal{O}_U) \ar[ld]^{\pi_U} \\
& (\Sh(\mathcal{D}/V), \mathcal{O}_V)
}
$$
Then there exists a morphism of topoi
$$
\sigma' : \Sh(\mathcal{D}/V) \to \Sh(\mathcal{C}'/U'),
$$
such that setting $\sigma = g' \circ \sigma'$ we have
$\pi'_{U'} \circ \sigma' = \text{id}$, $\pi_U \circ \sigma = \text{id}$,
$(\sigma')^{-1} = \pi'_{U', *}$, and $\sigma^{-1} = \pi_{U, *}$.
\end{lemma}

\begin{proof}
Let $v' : \mathcal{D}/V \to \mathcal{C}'/U'$ be the functor constructed
in the proof of Lemma \ref{lemma-fibred-category-with-object} starting
with $p' : \mathcal{C}' \to \mathcal{D}'$ and the object $U'$.
Since $u$ is a $1$-morphism of fibred categories over $\mathcal{D}$
it transforms strongly cartesian morphisms into strongly cartesian morphisms,
hence the functor $v = u \circ v'$ is the functor of
the proof of Lemma \ref{lemma-fibred-category-with-object}
relative to $p : \mathcal{C} \to \mathcal{D}$ and $U$. Thus our lemma
follows from that lemma.
\end{proof}

\end{document}
